\section{Coverage and applicability}\label{supp:coverage}

The benchmark scores methods on a six-game core drawn from a larger bit-exact substrate. The
coverage table makes the scope of that choice explicit. Paper~1 validates the
differentiable VCS ports bit-for-bit against the \texttt{xitari} reference on all 64 Arcade
Learning Environment games, within a 30-frame RAM and 60-frame screen window. The
interpretability study reported here runs on the six games listed in
Table~\ref{supp:tab:coverage}. Every method is scored on a shared gameplay state, produced by
one seeded random-action stream (seed~$0$, $90$-frame prefix, $15$-frame analysis horizon), not
on the boot screen. This keeps the analysis inside the Paper~1 bit-exact window while giving the
oracle a rich, non-degenerate cause set to score against. We pick these six games because all
three tiers of ground truth (\Sectionname\ref{sec:tiers}) are available for them and the oracle is
exact. The remaining 58 games are
bit-exact substrate but are not scored here. Some fall outside the validated horizon for the
analyses we run. For others, the tiered labels they would need are not yet produced. They are
candidates for the larger sweep, and are not excluded for any methodological reason. The
horizon column records the validated window. The tier column records which of the three
ground-truth tiers (T1 exact intervention, T2 data-flow, T3 externally-supplied semantic
labels) is available.

\begin{table}[t]
\centering
\caption{\textbf{Core-game coverage.} The six core games carry exact tiered ground truth and
the full method battery. Faithfulness is scored within the Paper~1 bit-exact window. The
last two columns give the methods run and the per-game output records contributed. The wider
64-game substrate is bit-exact (Paper~1) but is reserved for the larger sweep. T3 labels are
external (\Sectionname\ref{sec:tiers}); we record this so the benchmark does not appear to certify them.}
\label{supp:tab:coverage}
\footnotesize
\setlength{\tabcolsep}{4pt}
\begin{tabular}{@{}p{2.4cm}p{1.9cm}p{1.5cm}p{1.7cm}p{1.5cm}p{1.4cm}@{}}
\toprule
Game & Paper-1 status & Horizon & Tiers & Methods run & Records \\
\midrule
breakout & bit-exact & 30/60 fr & T1,T2,T3 & 30 + oracle & per-game \\
ms\_pacman & bit-exact & 30/60 fr & T1,T2,T3 & 30 + oracle & per-game \\
pong & bit-exact & 30/60 fr & T1,T2,T3 & 30 + oracle & per-game \\
qbert & bit-exact & 30/60 fr & T1,T2,T3 & 30 + oracle & per-game \\
seaquest & bit-exact & 30/60 fr & T1,T2,T3 & 30 + oracle & per-game \\
space\_invaders & bit-exact & 30/60 fr & T1,T2,T3 & 30 + oracle & per-game \\
\midrule
58 further games & bit-exact (P1) & 30/60 fr & T1,T2 partial & reserved & --- \\
\bottomrule
\end{tabular}
\end{table}

A gameplay state can still be degenerate for a particular method, and the benchmark records
this rather than scoring it as a failure. The oracle cause-density gate is the first guard. It
counts how many candidate causes move the shared output above a floor, and rejects a state that
does not clear a minimum count, so a near-flat oracle column never reaches a method. On
\texttt{seaquest} the gate raises the count of live candidate causes at the analysis frame from
$2$ to $19$, which is what makes the game usable. Three residual degenerate cells survive the
gate and are marked, not scored as zero. On \texttt{ms\_pacman} the pairwise neuroscience
analysis (A4) finds no coupled candidate pair at that frame, so its cell is uninformative. The
on-distribution counterfactual returns a zero delta on \texttt{space\_invaders},
\texttt{ms\_pacman}, and \texttt{qbert}, where its perturbation does not move the shared output;
those cells are recorded as not-valid. Expected Gradients scores low on the content regime of
every game but \texttt{pong}, because the causal byte it should attribute to is static there.
None of these cells is read as a wrong answer.

Whether a method applies at all depends on what the method needs from its target. The
applicability matrix records that for every method family. The Phase~A neuroscience analyses
read the raw machine state and need neither a learned policy nor gradients. The Phase~B
attribution methods split into two groups. The gradient members need a differentiable forward
pass. The intervention members need only the ability to set state. The Phase~C mechanistic
methods need access to internal activations. The causal members among them also need the
ability to intervene. Table~\ref{supp:tab:applic} records, per family, whether the method
operates on the raw VCS, whether it requires a learned policy, gradients, internal layers or
attention, or external labels, whether it uses oracle-like interventions, and whether it
receives a supplied hypothesis. The table shows the same pattern the results confirm. The
families that score high on faithfulness are the ones that can intervene. The families that
collapse on the position regime are the gradient and correlational ones.

\begin{table}[t]
\centering
\caption{\textbf{Per-family applicability.} What each method family requires of its target.
``Intervenes'' marks oracle-like $\mathrm{do}(u)$ access; ``hypothesis'' marks methods that
receive a supplied circuit hypothesis to test. The raw-VCS column distinguishes analyses that
read machine state directly from those that need a differentiable or layered model. A blank
means the requirement does not apply.}
\label{supp:tab:applic}
\footnotesize
\begin{tabular}{@{}p{3.0cm}cccccc@{}}
\toprule
Family & raw VCS & policy & gradients & layers/attn & intervenes & hypothesis \\
\midrule
A neuroscience & yes & no & no & no & A2 only & no \\
B gradient & yes & no & yes & no & no & no \\
B intervention & yes & no & no & no & yes & no \\
C causal & yes & no & no & yes & yes & yes \\
C gradient (attr.\ patch) & yes & no & yes & yes & approx. & yes \\
C probing & yes & no & no & yes & no & no \\
C / A dim.\ reduction & yes & no & no & yes & no & no \\
oracle (control) & yes & no & no & yes & exact & --- \\
\bottomrule
\end{tabular}
\end{table}
